%=======================================================================
%  Apart Research — Hackathon Submission Template (LaTeX port)
%  AI Incident Response Sprint, September 2026
%
%  Toggle \guidancetrue / \guidancefalse below to show or hide every
%  piece of italic guidance text at once. Set it to false before
%  submitting — the original template asks you to delete all guidance.
%=======================================================================

\documentclass[11pt,a4paper]{article}   % use [letterpaper] to match the original PDF

%------------------------- packages ------------------------------------
\usepackage[utf8]{inputenc}
\usepackage[T1]{fontenc}
\usepackage[english]{babel}
\IfFileExists{lmodern.sty}{\usepackage{lmodern}}{}  % optional: nicer fonts if installed
\usepackage[margin=1in]{geometry}
\usepackage{graphicx}
\usepackage{booktabs}
\usepackage{amsmath,amssymb}
\usepackage{enumitem}
\usepackage{xcolor}
\usepackage{caption}
\usepackage{titlesec}
\usepackage[hidelinks,breaklinks]{hyperref}

%------------------------- guidance switch -----------------------------
\newif\ifguidance
\guidancetrue          % <-- set to \guidancefalse for the final PDF

\definecolor{guidegray}{gray}{0.40}

% \guide{...} : inline italic guidance, disappears when guidance is off
\newcommand{\guide}[1]{\ifguidance\textcolor{guidegray}{\itshape #1}\par\medskip\fi}

% framed "how to use this template" box (vanishes entirely when guidance is off)
\usepackage{comment}
\ifguidance
  \newsavebox{\gbox}
  \newenvironment{guidancebox}
    {\begin{lrbox}{\gbox}\begin{minipage}{0.88\textwidth}%
     \small\color{guidegray}\itshape}
    {\end{minipage}\end{lrbox}%
     \begin{center}\setlength{\fboxsep}{10pt}%
     \fcolorbox{guidegray}{gray!6}{\usebox{\gbox}}\end{center}\medskip}
\else
  \excludecomment{guidancebox}
\fi

%------------------------- author block --------------------------------
% One entry per author. Add or remove \authorentry lines freely.
\newcommand{\authorentry}[2]{%
  \begin{minipage}[t]{0.30\linewidth}\centering
    #1\\[2pt]{\small\itshape #2}
  \end{minipage}\hspace{0.01\linewidth}}

%------------------------- section styling -----------------------------
\titleformat{\section}{\normalfont\large\bfseries}{\thesection.}{0.5em}{}
\titleformat{\subsection}{\normalfont\normalsize\bfseries}{\thesubsection}{0.5em}{}
\setlength{\parskip}{0.5em}
\setlength{\parindent}{0pt}

%=======================================================================
\begin{document}

%------------------------- title page ----------------------------------
\begin{center}
  {\LARGE\bfseries PROJECT TITLE\footnote{Research conducted at the AI Incident
   Response Sprint, September 2026.}}\\[1.5em]

  \authorentry{Author name 1}{Affiliation}
  \authorentry{Author name 2}{Affiliation}
  \authorentry{Author name 3}{Affiliation}\\[1.2em]
  \authorentry{Author name 4}{Affiliation}
  \authorentry{Author name 5}{Affiliation}
  \authorentry{Author name 6}{Affiliation}\\[1.5em]

  {\itshape With}\\[2pt]
  Apart Research
\end{center}

\vspace{1em}

%------------------------- abstract ------------------------------------
\begin{center}\textbf{Abstract}\end{center}
\begin{quote}
\guide{Summarize your project in 150--250 words. A strong abstract lets a reviewer
understand what you did and why it matters without reading anything else. Make sure
to cover: the problem, your approach, key results, and the main takeaway. Polish it
last: the abstract should reflect your final results, not your initial plan.}
\ifguidance\else\textit{[Abstract here.]}\fi
\end{quote}

\vspace{0.5em}

\begin{guidancebox}
\textbf{How to use this template:} Replace the italicized guidance text under each
section with your content. The section structure is strong guidance but not rigid.
If your project requires a different organization, feel free to adapt. Delete all
guidance text including this info box before submitting (set
\texttt{\textbackslash guidancefalse} at the top of the preamble). Make sure the
project title and author information above are up to date.

\medskip
\textbf{How your submission will be evaluated:} Your project will be judged on the
quality of this written report. To understand what we look for, see the Evaluation
Rubric.

\medskip
\textbf{Length:} the sprint page asks for 4 to 8 pages (8 max) excluding
references and appendices. The original template's rough guide for 4 pages:
Intro \& Related Work 1p, Methods and Results 2.5p, Discussion 0.5p.
\end{guidancebox}

%------------------------- 1. Introduction -----------------------------
\section{Introduction}
\guide{What problem are you addressing and why does it matter? Connect it to your
work: we want to know why your work is practically valuable. Provide enough
background for readers to understand your work. If relevant, briefly describe the
threat model or failure mode you're addressing; reference prior work that motivates
why it is worth addressing, or explain it yourself. Aspire to clearly list your most
important contributions that go beyond what exists today.}

Our main contributions are:
\begin{enumerate}[leftmargin=*,itemsep=2pt]
  \item {[First contribution --- what new thing did you create, discover, or demonstrate?]}
  \item {[Second contribution]}
  \item {[Third contribution, if applicable]}
\end{enumerate}

%------------------------- 2. Related Work -----------------------------
\section{Related Work}
\guide{What prior work is most similar, and how does your work differ? Cite the most
relevant papers, tools, or projects. Explain what gap your work addresses. Some
questions which may help: (a) When and why would someone use your method over the
existing state-of-the-art? (b) What information/insight does your method provide
which we did not have before?}

%------------------------- 3. Methods ----------------------------------
\section{Methods}
\guide{Describe your approach clearly enough that someone could replicate it.
Include key design choices and justify them where relevant (hint: the more you can
back up your design choices by referencing prior work, the better). What
models/datasets/tools did you use and why? What were key parameters or design
decisions? What did you try that didn't work? Could someone reproduce your work from
this description?}

%------------------------- 4. Results ----------------------------------
\section{Results}
\guide{Present your main findings with appropriate evidence. Use figures and tables
where appropriate (we strongly encourage at least one figure). Distinguish between
observations and interpretations. Argue why your claims are robust. E.g., if your
approach ``performs better'' than alternatives: (a) Do you have enough data? Is the
difference statistically significant? (b) Is it robust? Or do small changes to your
setup cause substantial changes to your results?}

\begin{guidancebox}
\textbf{Tips for figures and tables:} Number all figures (Figure 1, Figure 2\dots)
and tables (Table 1, Table 2\dots). Include descriptive captions that can be
understood without the main text. Place figures/tables near where they're first
referenced. Ensure text in figures is legible.
\end{guidancebox}

% ---- example figure / table scaffolding (delete if unused) ----
% \begin{figure}[t]
%   \centering
%   \includegraphics[width=0.8\linewidth]{figures/example.pdf}
%   \caption{Self-contained caption describing what the figure shows.}
%   \label{fig:example}
% \end{figure}
%
% \begin{table}[t]
%   \centering
%   \begin{tabular}{lrr}
%     \toprule
%     Model & Metric A & Metric B \\
%     \midrule
%     \dots & \dots & \dots \\
%     \bottomrule
%   \end{tabular}
%   \caption{Self-contained caption describing what the table shows.}
%   \label{tab:example}
% \end{table}

%------------------------- 5. Discussion -------------------------------
\section{Discussion and Limitations}
\guide{Discuss the broader implications for AI safety. What do your results mean?
What trends do you notice and what might they indicate?}

\subsection*{Limitations}
\guide{What are the limitations of your work? What threat models or failure modes
did you not address? Be honest about constraints --- methodological limitations,
scope limitations, or aspects you couldn't fully address in the hackathon timeframe.
Explicitly note the assumptions you made, whether implicitly or explicitly, and how
the interpretation of your results would change if a given assumption did not hold.}

\subsection*{Future Work}
\guide{What are the natural next steps? How could this work be extended?}

%------------------------- 6. Conclusion -------------------------------
\section{Conclusion}
\guide{Briefly summarize your main findings and their implications (1--2 paragraphs).}

%------------------------- Code and Data -------------------------------
\section*{Code and Data}
\guide{Include links if applicable. If your project doesn't involve code (e.g.,
policy analysis) or if there are info-hazard considerations, note that here.}
\begin{itemize}[leftmargin=*,itemsep=2pt]
  \item Code repository: [Link to GitHub/GitLab if applicable]
  \item Data/Datasets: [Link if applicable]
  \item Other artifacts (optional): [Demo link, video walkthrough, Hugging Face Space, etc.]
\end{itemize}

%------------------------- Author Contributions ------------------------
\section*{Author Contributions (optional)}
\guide{e.g., ``A.B. led the project and designed experiments. C.D. implemented the
code. All authors contributed to writing and reviewed the final manuscript.''}

%------------------------- References ----------------------------------
\section*{References}
\guide{Use a consistent citation format. Include: Author(s), Year, Title,
Venue/Publisher, and URL or DOI where available.}

\begin{enumerate}[leftmargin=*,itemsep=2pt]
  \item {[Reference 1]}
  \item {[Reference 2]}
  \item \dots
\end{enumerate}

% If you prefer BibTeX, delete the block above and uncomment:
% \bibliographystyle{plainnat}
% \bibliography{refs}

%------------------------- Appendix ------------------------------------
\appendix
\section*{Appendix (optional)}
\guide{Supplementary material such as additional figures, detailed methodology,
prompts used, extended results, etc.}

%------------------------- Dual-use appendix ---------------------------
% NOTE: the sprint's submission requirements list a mandatory
% "Limitations and Dual-Use Considerations" appendix, which the official
% template does not include as its own section. Keeping it explicit here.
\section*{Limitations and Dual-Use Considerations}
\guide{Required by the sprint submission checklist. State what your artifact does
and does not establish, and any way it could be misused. If you withheld results
(e.g.\ novel installation recipes) pending disclosure review, say so here.}

%------------------------- LLM usage -----------------------------------
\section*{LLM Usage Statement}
\guide{If you used LLM assistance in developing your project or writing this report,
briefly note how. Ensure all claims and results have been verified. NOTE: We
strongly encourage that the final version of the submission is primarily written by
your team. E.g., ``We used Claude to brainstorm approaches and help draft sections.
All results and claims were independently verified.''}

\end{document}
